% Canonical LaTeX source; imported and revised from the legacy Markdown.
\chapter{Primitive e integrali indefiniti}
\label{chap:14}

\begin{obiettivocapitolo}{riconoscere la struttura di un integrale indefinito e verificarne la primitiva}
\prioritaalta. Non esiste un algoritmo unico: la scelta del metodo nasce dalla forma dell'integrando.
\end{obiettivocapitolo}
\begin{sceltametodo}{quale tecnica tentare}
\begin{longtable}{@{}p{.29\textwidth}p{.31\textwidth}p{.31\textwidth}@{}}
\toprule
\textbf{Struttura} & \textbf{Metodo} & \textbf{Segnale riconoscibile}\\
\midrule\endhead
Tabella elementare & primitiva immediata & potenze, esponenziali, seno/coseno standard\\
$f'(x)\,g(f(x))$ & sostituzione & derivata della funzione interna presente a fattore\\
Prodotto con logaritmo, polinomio, inversa trigonometrica & parti & una derivazione semplifica un fattore\\
Razionale $P/Q$ & divisione e fratti semplici & grado, fattorizzazione reale del denominatore\\
Potenze di seno e coseno & identità e riduzioni & parità degli esponenti\\
Radici di quadratiche & sostituzione trigonometrica o iperbolica & forme $a^2-x^2$, $a^2+x^2$, $x^2-a^2$\\
\bottomrule
\end{longtable}
\end{sceltametodo}
\begin{proceduraesame}{integrazione per sostituzione}
Scegli $u=\varphi(x)$, calcola $du=\varphi'(x)dx$, riscrivi integralmente l'integrando in $u$, integra e torna a $x$. Se restano contemporaneamente $u$ e $x$, la sostituzione non è stata completata.
\end{proceduraesame}
\begin{proceduraesame}{funzione razionale}
Esegui la divisione se $\deg P\ge\deg Q$; fattorizza $Q$ sui reali; decomponi con un termine per ogni potenza dei fattori lineari e con numeratore lineare per i quadratici irriducibili; determina i coefficienti e integra termine a termine.
\end{proceduraesame}
\begin{controllofinale}{primitiva}
Deriva il risultato e recupera l'integrando sul dominio considerato. Ricorda la costante $C$ e le possibili costanti differenti sulle componenti connesse.
\end{controllofinale}

\section{Primitive e costanti}\label{primitive-e-costanti}

Su un intervallo \(I\):

\[
\displaystyle
F'(x)=f(x)
\quad\Longleftrightarrow\quad
\int f(x)\,dx=F(x)+C.
\]

Equivalentemente:

\[
\frac{d}{dx}\left[\int f(x)\,dx\right]=f(x),
\qquad
\int F'(x)\,dx=F(x)+C,
\]

intendendo il primo membro come una qualsiasi primitiva scelta localmente.

Se \(D=\bigcup_j I_j\) è unione disgiunta di componenti intervallari:

\[
\displaystyle
G_{\mid I_j}=F_{\mid I_j}+C_j,
\qquad C_j\in\mathbb R.
\]

Linearità. Se \(F'=f\) e \(G'=g\), allora per ogni \(\alpha,\beta\in\mathbb R\):

\[
\displaystyle
\int\bigl(\alpha f(x)+\beta g(x)\bigr)\,dx
=
\alpha F(x)+\beta G(x)+C.
\]

Teoria: \href{https://ingegnerismo.it/matematica/integrazione/}{integrazione} e \href{https://ingegnerismo.it/matematica/integrali-immediati/}{integrali immediati}.

\section{Tabella delle primitive fondamentali}\label{tabella-delle-primitive-fondamentali}

Nelle formule seguenti \(a\ne0\); le costanti additive sono intese per ogni componente del dominio.

\begin{longtable}[]{@{}
  >{\raggedright\arraybackslash}p{(\columnwidth - 4\tabcolsep) * \real{0.3333}}
  >{\raggedright\arraybackslash}p{(\columnwidth - 4\tabcolsep) * \real{0.3333}}
  >{\raggedright\arraybackslash}p{(\columnwidth - 4\tabcolsep) * \real{0.3333}}@{}}
\toprule\noalign{}
\begin{minipage}[b]{\linewidth}\raggedright
Integranda
\end{minipage} & \begin{minipage}[b]{\linewidth}\raggedright
Primitiva
\end{minipage} & \begin{minipage}[b]{\linewidth}\raggedright
Condizioni
\end{minipage} \\
\midrule\noalign{}
\endhead
\bottomrule\noalign{}
\endlastfoot
\(1\) & \(x+C\) & \(x\in\mathbb R\) \\
\(x^\alpha\) & \(\dfrac{x^{\alpha+1}}{\alpha+1}+C\) & \(\alpha\ne-1\); sugli intervalli reali di definizione \\
\(\dfrac1x\) & \(\ln\lvert x\rvert+C\) & \(x\ne0\) \\
\((ax+b)^\alpha\) & \(\dfrac{(ax+b)^{\alpha+1}}{a(\alpha+1)}+C\) & \(\alpha\ne-1\); sugli intervalli di definizione \\
\(\dfrac1{ax+b}\) & \(\dfrac1a\ln\lvert ax+b\rvert+C\) & \(ax+b\ne0\) \\
\(e^{ax+b}\) & \(\dfrac1a e^{ax+b}+C\) & \(x\in\mathbb R\) \\
\(q^{ax+b}\) & \(\dfrac{q^{ax+b}}{a\ln q}+C\) & \(q>0,\ q\ne1\) \\
\(\sin(ax+b)\) & \(-\dfrac1a\cos(ax+b)+C\) & \(x\in\mathbb R\) \\
\(\cos(ax+b)\) & \(\dfrac1a\sin(ax+b)+C\) & \(x\in\mathbb R\) \\
\(\tan(ax+b)\) & \(-\dfrac1a\ln\lvert\cos(ax+b)\rvert+C\) & \(\cos(ax+b)\ne0\) \\
\(\cot(ax+b)\) & \(\dfrac1a\ln\lvert\sin(ax+b)\rvert+C\) & \(\sin(ax+b)\ne0\) \\
\(\sec(ax+b)\) & \(\dfrac1a\ln\lvert\sec(ax+b)+\tan(ax+b)\rvert+C\) & \(\cos(ax+b)\ne0\) \\
\(\csc(ax+b)\) & \(\dfrac1a\ln\lvert\csc(ax+b)-\cot(ax+b)\rvert+C\) & \(\sin(ax+b)\ne0\) \\
\(\sec^2(ax+b)\) & \(\dfrac1a\tan(ax+b)+C\) & \(\cos(ax+b)\ne0\) \\
\(\csc^2(ax+b)\) & \(-\dfrac1a\cot(ax+b)+C\) & \(\sin(ax+b)\ne0\) \\
\(\sec(ax+b)\tan(ax+b)\) & \(\dfrac1a\sec(ax+b)+C\) & \(\cos(ax+b)\ne0\) \\
\(\csc(ax+b)\cot(ax+b)\) & \(-\dfrac1a\csc(ax+b)+C\) & \(\sin(ax+b)\ne0\) \\
\(\sinh(ax+b)\) & \(\dfrac1a\cosh(ax+b)+C\) & \(x\in\mathbb R\) \\
\(\cosh(ax+b)\) & \(\dfrac1a\sinh(ax+b)+C\) & \(x\in\mathbb R\) \\
\(\tanh(ax+b)\) & \(\dfrac1a\ln\cosh(ax+b)+C\) & \(x\in\mathbb R\) \\
\(\coth(ax+b)\) & \(\dfrac1a\ln\lvert\sinh(ax+b)\rvert+C\) & \(\sinh(ax+b)\ne0\) \\
\end{longtable}

Per \(a>0\):

\begin{longtable}[]{@{}
  >{\raggedright\arraybackslash}p{(\columnwidth - 4\tabcolsep) * \real{0.3333}}
  >{\raggedright\arraybackslash}p{(\columnwidth - 4\tabcolsep) * \real{0.3333}}
  >{\raggedright\arraybackslash}p{(\columnwidth - 4\tabcolsep) * \real{0.3333}}@{}}
\toprule\noalign{}
\begin{minipage}[b]{\linewidth}\raggedright
Integranda
\end{minipage} & \begin{minipage}[b]{\linewidth}\raggedright
Primitiva
\end{minipage} & \begin{minipage}[b]{\linewidth}\raggedright
Condizioni
\end{minipage} \\
\midrule\noalign{}
\endhead
\bottomrule\noalign{}
\endlastfoot
\(\dfrac1{a^2+x^2}\) & \(\dfrac1a\arctan\dfrac xa+C\) & \(x\in\mathbb R\) \\
\(\dfrac1{\sqrt{a^2-x^2}}\) & \(\arcsin\dfrac xa+C\) & \(\lvert x\rvert<a\) \\
\(\dfrac1{a^2-x^2}\) & \(\dfrac1{2a}\ln\left\lvert\dfrac{a+x}{a-x}\right\rvert+C\) & \(x\ne\pm a\) \\
\(\dfrac1{x^2-a^2}\) & \(\dfrac1{2a}\ln\left\lvert\dfrac{x-a}{x+a}\right\rvert+C\) & \(x\ne\pm a\) \\
\(\dfrac1{\sqrt{x^2+a^2}}\) & \(\operatorname{arsinh}\dfrac xa+C=\ln\left\lvert x+\sqrt{x^2+a^2}\right\rvert+C\) & \(x\in\mathbb R\) \\
\(\dfrac1{\sqrt{x^2-a^2}}\) & \(\ln\left\lvert x+\sqrt{x^2-a^2}\right\rvert+C\) & \(\lvert x\rvert>a\) \\
\end{longtable}

Teoria e tabelle estese: \href{https://ingegnerismo.it/matematica/integrali-immediati/}{integrali immediati}.

\section{Primitive di funzioni composte}\label{primitive-di-funzioni-composte}

Schemi immediati ottenuti con la regola della catena inversa:

\begin{longtable}[]{@{}
  >{\raggedright\arraybackslash}p{(\columnwidth - 4\tabcolsep) * \real{0.3333}}
  >{\raggedright\arraybackslash}p{(\columnwidth - 4\tabcolsep) * \real{0.3333}}
  >{\raggedright\arraybackslash}p{(\columnwidth - 4\tabcolsep) * \real{0.3333}}@{}}
\toprule\noalign{}
\begin{minipage}[b]{\linewidth}\raggedright
Integranda
\end{minipage} & \begin{minipage}[b]{\linewidth}\raggedright
Primitiva
\end{minipage} & \begin{minipage}[b]{\linewidth}\raggedright
Condizioni
\end{minipage} \\
\midrule\noalign{}
\endhead
\bottomrule\noalign{}
\endlastfoot
\(u'(x)u(x)^\alpha\) & \(\dfrac{u(x)^{\alpha+1}}{\alpha+1}+C\) & \(\alpha\ne-1\); dominio reale della potenza \\
\(\dfrac{u'(x)}{u(x)}\) & \(\ln\lvert u(x)\rvert+C\) & \(u(x)\ne0\) \\
\(u'(x)e^{u(x)}\) & \(e^{u(x)}+C\) & --- \\
\(u'(x)a^{u(x)}\) & \(\dfrac{a^{u(x)}}{\ln a}+C\) & \(a>0\), \(a\ne1\) \\
\(u'(x)\cos u(x)\) & \(\sin u(x)+C\) & --- \\
\(u'(x)\sin u(x)\) & \(-\cos u(x)+C\) & --- \\
\(\dfrac{u'(x)}{1+u(x)^2}\) & \(\arctan u(x)+C\) & --- \\
\(\dfrac{u'(x)}{\sqrt{1-u(x)^2}}\) & \(\arcsin u(x)+C\) & \(\lvert u(x)\rvert<1\) \\
\(-\dfrac{u'(x)}{\sqrt{1-u(x)^2}}\) & \(\arccos u(x)+C\) & \(\lvert u(x)\rvert<1\) \\
\end{longtable}

\section{Sostituzione}\label{sostituzione}

Se \(g\in C^1(I)\) e \(f\) è continua su un intervallo contenente \(g(I)\):

\[
\displaystyle
u=g(x),
\qquad
\int f(g(x))g'(x)\,dx
=
\int f(u)\,du.
\]

Se \(g\in C^1([a,b])\) e \(f\) è continua su un intervallo contenente \(g([a,b])\):

\[
\displaystyle
\int_a^b f(g(x))g'(x)\,dx
=
\int_{g(a)}^{g(b)}f(u)\,du.
\]

Con \(x=\varphi(t)\), \(\varphi\in C^1\), se \(F'=f\) allora

\[
\displaystyle
\int f(\varphi(t))\varphi'(t)\,dt
=
F(\varphi(t))+C.
\]

Se \(\varphi\) è invertibile, al termine si sostituisce \(t=\varphi^{-1}(x)\).

Teoria ed esempi: \href{https://ingegnerismo.it/matematica/integrazione-per-sostituzione/}{integrazione per sostituzione}.

\section{Integrazione per parti}\label{integrazione-per-parti}

Per \(u,v\in C^1(I)\):

\[
\displaystyle
\int u(x)v'(x)\,dx
=
u(x)v(x)-\int u'(x)v(x)\,dx,
\]

\[
\displaystyle
\int u\,dv=uv-\int v\,du.
\]

Forma definita:

\[
\displaystyle
\int_a^b u(x)v'(x)\,dx
=
\bigl[u(x)v(x)\bigr]_a^b
-\int_a^b u'(x)v(x)\,dx.
\]

Teoria, scelta dei fattori e casi ripetuti: \href{https://ingegnerismo.it/matematica/integrazione-per-parti/}{integrazione per parti}.

\section{Funzioni razionali e fratti semplici}\label{funzioni-razionali-e-fratti-semplici}

Per \(\deg P\ge\deg Q\):

\[
\displaystyle
\frac{P(x)}{Q(x)}
=
S(x)+\frac{R(x)}{Q(x)},
\qquad
\deg R<\deg Q.
\]

Se

\[
\displaystyle
Q(x)
=
c_0\prod_i(x-a_i)^{m_i}
\prod_j(x^2+p_jx+r_j)^{n_j},
\qquad
c_0\ne0,
\qquad
p_j^2-4r_j<0,
\]

allora

\[
\displaystyle
\frac{R(x)}{Q(x)}
=
\sum_i\sum_{k=1}^{m_i}\frac{A_{ik}}{(x-a_i)^k}
+
\sum_j\sum_{k=1}^{n_j}
\frac{B_{jk}x+C_{jk}}{(x^2+p_jx+r_j)^k}.
\]

Primitive dei fattori lineari:

\[
\displaystyle
\int\frac{A}{x-a}\,dx
=
A\ln\lvert x-a\rvert+C,
\]

\[
\displaystyle
\int\frac{A}{(x-a)^k}\,dx
=
\frac{A}{1-k}(x-a)^{1-k}+C,
\qquad k\ge2.
\]

Per \(Q_2(x)=x^2+px+r\), \(p^2-4r<0\):

\[
\displaystyle
Bx+C
=
\frac B2(2x+p)+\left(C-\frac{Bp}{2}\right).
\]

La prima parte è proporzionale a \(Q_2'(x)\) e produce un logaritmo; la parte costante residua produce un'arcotangente dopo il completamento del quadrato.

Teoria e determinazione dei coefficienti: \href{https://ingegnerismo.it/matematica/fratti-semplici/}{fratti semplici}.

\section{Primitive dei fattori quadratici irriducibili}\label{primitive-dei-fattori-quadratici-irriducibili}

Per evitare di usare la stessa lettera per il polinomio e per il termine noto, poniamo

\[
Q_2(x)=x^2+px+r,
\qquad
p^2-4r<0,
\qquad
h=\sqrt{r-\frac{p^2}{4}}>0.
\]

Allora

\[
\int\frac{Bx+C}{Q_2(x)}\,dx
=
\frac B2\ln Q_2(x)
+
\frac{C-Bp/2}{h}
\arctan\left(\frac{x+p/2}{h}\right)+C_0.
\]

Qui \(B,C\) sono coefficienti del numeratore, mentre \(C_0\) è la costante di integrazione.

\section{Integrali trigonometrici e formule di riduzione}\label{integrali-trigonometrici-e-formule-di-riduzione}

Per \(n\ge2\), formule di riduzione:

\[
\int\sin^n x\,dx
=-\frac{\sin^{n-1}x\cos x}{n}
+\frac{n-1}{n}\int\sin^{n-2}x\,dx,
\]

\[
\int\cos^n x\,dx
=\frac{\cos^{n-1}x\sin x}{n}
+\frac{n-1}{n}\int\cos^{n-2}x\,dx.
\]

Per \(\int\sin^m x\cos^n x\,dx\):

\begin{itemize}
\tightlist
\item
  se uno degli esponenti è dispari, si conserva un fattore della funzione corrispondente e si usa \(\sin^2x=1-\cos^2x\) oppure \(\cos^2x=1-\sin^2x\);
\item
  se entrambi sono pari, si usano le formule di bisezione.
\end{itemize}

\section{Sostituzioni trigonometriche}\label{sostituzioni-trigonometriche}

Sostituzione di Weierstrass:

\[
\displaystyle
t=\tan\frac x2,
\qquad
\sin x=\frac{2t}{1+t^2},
\qquad
\cos x=\frac{1-t^2}{1+t^2},
\qquad
dx=\frac{2}{1+t^2}\,dt,
\]

\[
\displaystyle
x\ne(2k+1)\pi,
\qquad k\in\mathbb Z.
\]

Per \(a>0\):

\begin{longtable}[]{@{}
  >{\raggedright\arraybackslash}p{(\columnwidth - 4\tabcolsep) * \real{0.3333}}
  >{\raggedright\arraybackslash}p{(\columnwidth - 4\tabcolsep) * \real{0.3333}}
  >{\raggedright\arraybackslash}p{(\columnwidth - 4\tabcolsep) * \real{0.3333}}@{}}
\toprule\noalign{}
\begin{minipage}[b]{\linewidth}\raggedright
Radicale
\end{minipage} & \begin{minipage}[b]{\linewidth}\raggedright
Sostituzione
\end{minipage} & \begin{minipage}[b]{\linewidth}\raggedright
Identità risultante
\end{minipage} \\
\midrule\noalign{}
\endhead
\bottomrule\noalign{}
\endlastfoot
\(\sqrt{a^2-x^2}\) & \(x=a\sin\theta\), \(-\dfrac\pi2\le\theta\le\dfrac\pi2\) & \(\sqrt{a^2-x^2}=a\cos\theta\), \(dx=a\cos\theta\,d\theta\) \\
\(\sqrt{a^2+x^2}\) & \(x=a\tan\theta\), \(-\dfrac\pi2<\theta<\dfrac\pi2\) & \(\sqrt{a^2+x^2}=a\sec\theta\), \(dx=a\sec^2\theta\,d\theta\) \\
\(\sqrt{x^2-a^2}\), \(x\ge a\) & \(x=a\sec\theta\), \(0\le\theta<\dfrac\pi2\) & \(\sqrt{x^2-a^2}=a\tan\theta\), \(dx=a\sec\theta\tan\theta\,d\theta\) \\
\(\sqrt{x^2-a^2}\), \(x\le-a\) & \(x=-a\sec\theta\), \(0\le\theta<\dfrac\pi2\) & \(\sqrt{x^2-a^2}=a\tan\theta\), \(dx=-a\sec\theta\tan\theta\,d\theta\) \\
\end{longtable}

Teoria e casi operativi: \href{https://ingegnerismo.it/matematica/sostituzioni-di-weierstrass/}{sostituzioni di Weierstrass} e \href{https://ingegnerismo.it/matematica/integrali-trigonometrici/}{integrali trigonometrici}.
